\section{System Architecture}\label{sec:architecture}

\prismloc{} is organized as two layers realized by four packages
(Fig.~\ref{fig:layers}). The lower layer holds the estimator cores:
\texttt{prism\_loc\_core}, which implements the Monte Carlo
localization (MCL) family, and \texttt{prism\_loc\_fusion}, which
implements the error-state Kalman filter (ESKF). Both are plain
C++17 libraries whose only third-party dependency is Eigen; each
\texttt{CMakeLists.txt} calls \texttt{find\_package(Eigen3 REQUIRED)}
and links \texttt{Eigen3::Eigen} alone, with no reference to
\texttt{rclcpp}, \texttt{tf2}, or PCL. The upper layer holds the ROS~2
node packages: \texttt{prism\_loc} (the \texttt{laser2d} and
\texttt{ndt3d} backends) and \texttt{prism\_loc\_fusion\_ros} (the
\texttt{fusion3d} backend). These packages, and only these, depend on
\texttt{rclcpp}, the \texttt{tf2} stack, the ROS message packages, and
PCL. Each core is consumed by its node package through the ordinary
ament \texttt{find\_package} mechanism, so the dependency direction is
strictly upward: the cores never link against the middleware.

\begin{figure}[t]
\centering
\footnotesize
\setlength{\tabcolsep}{4pt}
\begin{tabular}{|l|}
\hline
\textbf{Estimator cores} --- C++17 + Eigen, no ROS/PCL, gtest \\
\hline
\texttt{prism\_loc\_core} (MCL):\\
\quad OdometryMotionModel $\rightarrow$ ParticleFilter\\
\quad\quad $\vdash$ Laser2DLikelihoodField\\
\quad\quad $\vdash$ Ndt3DModel\\
\quad $\rightarrow$ Pose2D $+$ $6{\times}6$ covariance\\
\texttt{prism\_loc\_fusion} (ESKF):\\
\quad IMU predict $\rightarrow$ 15-state ESKF\\
\quad\quad $\dashv$ NDT 6-DoF pose update\\
\quad\quad $\dashv$ RTK-GNSS position update (WGS84 $\rightarrow$ ENU)\\
\hline
\multicolumn{1}{c}{}\\[-6pt]
\hline
\textbf{ROS~2 nodes} --- rclcpp / tf2 / PCL only here \\
\hline
\texttt{prism\_loc}\hfill \texttt{backend = laser2d | ndt3d}\\
\texttt{prism\_loc\_fusion\_ros}\hfill \texttt{backend = fusion3d}\\
\quad adapters: laser / pointcloud / ndt\_registration\\
\hline
in:~\texttt{/scan|/points|/imu|/gnss}, map, \texttt{/initialpose}, TF\\
out:~\texttt{/tf} (\texttt{map}$\rightarrow$\texttt{odom}), \texttt{\char`~/pose},
\texttt{\char`~/particle\_cloud|\char`~/odometry}\\
\hline
\end{tabular}
\caption{The two-layer split. Estimator cores (bottom) are
middleware-free libraries; the ROS~2 nodes (top) select a backend at
runtime and bridge messages to core types through adapters. The
dependency arrow points only upward.}
\label{fig:layers}
\end{figure}

\textbf{Why the split.} Isolating the estimators from the middleware
buys three properties. First, unit-testability without a running
graph: each core builds and tests with a plain system toolchain and
gtest, so its numerical behavior is exercised with no ROS installation,
no node, and no message plumbing. The \texttt{CMakeLists.txt} of each
core wires tests in dual mode --- through ament when a ROS workspace is
present, and through bare CTest plus system gtest otherwise --- and this
ROS-free path is the one continuous integration builds. The detailed
testing story is deferred to Section~\ref{sec:validation}. Second,
determinism: all sampling routes through a single seeded
\texttt{std::mt19937\_64} wrapper (default seed \texttt{42}), so a
motion-prediction, weighting, or resampling step is reproducible bit
for bit across runs, which is what makes the core tests assertable
rather than merely statistical. Third, one output contract shared by
all three backends. Regardless of backend, a node publishes the
\texttt{map}$\rightarrow$\texttt{odom} transform on \texttt{/tf}, a
\texttt{PoseWithCovarianceStamped} on \texttt{\char`~/pose}, and
accepts a \texttt{PoseWithCovarianceStamped} on \texttt{/initialpose}
for seeding; the MCL nodes additionally publish a particle cloud while
the fusion node publishes an \texttt{Odometry} message carrying pose
and velocity. Because the contract is fixed, choosing between the
likelihood-field MCL, NDT-MCL, and ESKF estimators is a runtime
decision: \texttt{prism\_loc} reads a \texttt{backend} string parameter
(\texttt{"laser2d"} by default) and branches on it when constructing
its measurement model and when applying the correction step, and
\texttt{prism\_loc\_fusion\_ros} supplies the \texttt{fusion3d}
estimator behind the identical topic and TF interface. As stated in
Section~\ref{sec:related}, this factoring --- three known estimator
families behind one ROS~2 contract --- is the architectural
contribution; the estimators themselves are textbook.

\textbf{Core interfaces.} The MCL core expresses the shared
predict--weight--resample loop against a small set of value types.
A \texttt{Pose2D} carries the planar pose $(x, y, \theta)$; a \texttt{Particle}
pairs a pose with a weight; and the two backends differ only through a
\texttt{MeasurementModel} abstraction whose single virtual method
returns the log-likelihood of an observation given a candidate pose.
\texttt{Laser2DLikelihoodField} (a 2D scan against an occupancy-grid
likelihood field) and \texttt{Ndt3DModel} (a 3D cloud against a
Gaussian-voxel NDT map) are its two implementations, and the same
\texttt{ParticleFilter} and \texttt{OdometryMotionModel} drive both.
The fusion core exposes an \texttt{Eskf} class over a 15-dimensional
error state with three entry points --- an IMU-driven
\texttt{predict}, a 6-DoF \texttt{updatePose} fed by NDT registration,
and a 3-DoF \texttt{updatePosition} fed by an absolute position fix ---
alongside a geodetic helper that converts WGS84 latitude/longitude/
altitude to a local ENU frame. Section~\ref{sec:cores} details these
estimators.

\textbf{Adapters.} The node layer bridges ROS message types to the
middleware-free core types through thin adapters, keeping the cores
ignorant of message definitions. In \texttt{prism\_loc}, a laser
adapter converts a \texttt{LaserScan} to the core's \texttt{LaserScan2D}
(honoring range overrides) and an \texttt{OccupancyGrid} to the core
grid map, and a pointcloud adapter converts a \texttt{PointCloud2}, and
a \texttt{.pcd} map file, into the \texttt{std::vector<Eigen::Vector3d>}
the NDT map ingests; both also convert quaternion-bearing
\texttt{Transform} and \texttt{Pose} messages to \texttt{Pose2D} by
extracting yaw. In \texttt{prism\_loc\_fusion\_ros}, an
\texttt{ndt\_registration} adapter wraps PCL's NDT to align an incoming
cloud against the prior-map cloud and returns the resulting
\texttt{Isometry3d} pose, which becomes the ESKF's 6-DoF measurement.
PCL therefore appears only inside these adapters, never in the cores.
The node-layer responsibilities --- TF lookups, topic wiring, and
runtime diagnostics --- are described in Section~\ref{sec:ros}.
